\chapter{METODOLOGIA}

Este capítulo apresenta os procedimentos metodológicos adotados para o desenvolvimento e a análise do simulador computacional. São descritos a caracterização da pesquisa, as etapas realizadas, a área de estudo e os dados utilizados, os procedimentos de modelagem e implementação, as estratégias de verificação e o planejamento dos experimentos.

A metodologia busca estabelecer uma separação entre o modelo conceitual, sua implementação computacional e os resultados obtidos durante as execuções. Essa distinção é importante para evitar que a existência de uma funcionalidade no código ou a produção de uma saída pelo simulador seja interpretada, isoladamente, como evidência de validação epidemiológica.

\section{Caracterização da pesquisa}

O presente trabalho caracteriza-se como uma pesquisa aplicada, de natureza computacional e caráter exploratório. O estudo envolve o desenvolvimento de um simulador espacial e estocástico baseado em agentes para representar aspectos da dinâmica de transmissão da dengue no município de São Paulo.

A pesquisa utiliza técnicas de modelagem e simulação computacional para representar entidades, estados, interações e processos temporais relacionados ao sistema estudado. A modelagem baseada em agentes foi adotada por permitir a representação de entidades heterogêneas e suas interações com outras entidades e com o ambiente.

O trabalho também possui caráter exploratório, uma vez que busca investigar o comportamento do modelo sob condições controladas e analisar as possibilidades e limitações decorrentes das hipóteses, dos dados e das regras adotadas.

A avaliação do simulador é tratada em diferentes níveis. Inicialmente, busca-se verificar a consistência entre as regras especificadas e a implementação computacional. Posteriormente, são realizados experimentos para analisar o comportamento produzido pelo modelo. A eventual comparação entre resultados simulados e dados epidemiológicos observados constitui uma etapa distinta e depende da compatibilidade entre indicadores, escalas espaciais e temporais e demais condições necessárias para essa análise.

\section{Etapas da pesquisa}

O desenvolvimento da pesquisa foi organizado em etapas relacionadas entre si, conforme apresentado a seguir:

\begin{enumerate}
    \item realização da revisão bibliográfica sobre dengue, modelagem e simulação computacional, modelagem baseada em agentes e representação espacial;
    
    \item definição do problema de pesquisa, dos objetivos e do escopo do estudo;
    
    \item obtenção e preparação dos dados geográficos e climáticos utilizados pelo simulador;
    
    \item especificação conceitual das entidades, estados, interações e processos que compõem o modelo;
    
    \item desenvolvimento e organização da implementação utilizando a plataforma GAMA e a linguagem GAML;
    
    \item definição de procedimentos de verificação para analisar regras, estados, parâmetros e saídas do simulador;
    
    \item planejamento e realização dos experimentos computacionais;
    
    \item coleta e análise dos indicadores produzidos pelas simulações;
    
    \item identificação das limitações do modelo e dos procedimentos adotados.
\end{enumerate}

As etapas não devem ser compreendidas como um processo estritamente linear. A construção de modelos computacionais pode exigir revisões nas regras, nos parâmetros e na implementação à medida que novas verificações são realizadas.

\section{Área de estudo e dados utilizados}

A área de estudo corresponde ao município de São Paulo, localizado no Estado de São Paulo, Brasil. A escolha do município está relacionada à dimensão territorial e populacional da cidade e à possibilidade de utilizar informações espaciais na representação do ambiente da simulação.

A representação computacional do território não busca reproduzir integralmente todas as características geográficas, demográficas e epidemiológicas do município. O espaço utilizado pelo simulador é definido a partir das bases geográficas selecionadas e das simplificações necessárias para atender aos objetivos do estudo.

\subsection{Dados geográficos}

Os dados geográficos são utilizados para representar o espaço da simulação e orientar a distribuição ou localização dos componentes do modelo. A preparação desses dados inclui a verificação das geometrias utilizadas e de sua compatibilidade com a área de estudo.

A utilização de uma base geográfica em uma simulação não implica, por si só, uma representação detalhada de todos os elementos presentes no território. A resolução espacial adotada depende das informações disponíveis e das decisões de modelagem.

Os procedimentos de preparação dos dados devem preservar a identificação da fonte utilizada e das transformações realizadas antes de sua incorporação ao simulador. Essas informações são necessárias para garantir a rastreabilidade entre os dados de entrada e a representação espacial adotada.

\subsection{Dados climáticos}

Os dados climáticos são utilizados como entradas para representar condições ambientais que podem influenciar processos relacionados ao vetor e ao ambiente da simulação.

Antes de sua utilização, os dados devem ser organizados em formato compatível com o processamento realizado pelo modelo. Esse procedimento pode incluir a seleção das variáveis de interesse, a padronização das datas e a verificação da consistência dos valores utilizados.

A utilização de dados climáticos históricos não significa que todas as relações entre clima e dinâmica epidemiológica sejam reproduzidas automaticamente. Os efeitos dessas variáveis dependem das regras implementadas e dos parâmetros definidos no modelo.

\section{Procedimento de modelagem}

O procedimento de modelagem consiste na definição da representação conceitual do sistema antes de sua descrição detalhada em código. Essa etapa estabelece quais entidades serão consideradas, quais atributos e estados podem possuir e quais regras representam suas interações.

A especificação do modelo é orientada pelos conceitos discutidos na fundamentação teórica e pelas informações disponíveis para o estudo. As entidades representam componentes considerados relevantes para o objetivo da simulação, enquanto os comportamentos representam abstrações dos processos que se pretende investigar.

A documentação do modelo utiliza os princípios do protocolo ODD como referência para a organização das informações. Dessa forma, são considerados aspectos como o propósito do modelo, as entidades, as escalas temporais e espaciais, os processos, as condições de inicialização, as entradas e os submodelos \cite{grimm_etal_2020}.

A modelagem também envolve a definição explícita das hipóteses e simplificações adotadas. Processos que não são representados diretamente ou que são descritos de forma agregada devem ser identificados para que os resultados sejam interpretados de acordo com o escopo do modelo.

A descrição detalhada das entidades, dos estados, das interações e das regras conceituais é apresentada no Capítulo 4.

\section{Procedimentos de implementação}

A implementação do simulador é realizada utilizando a plataforma GAMA e a linguagem GAML. O desenvolvimento é organizado em componentes destinados à representação dos agentes e dos processos globais da simulação.

Os procedimentos de implementação incluem a organização dos arquivos do projeto, a definição das espécies de agentes, a inicialização do ambiente e das populações e a implementação das regras responsáveis pela evolução da simulação.

Também são implementados mecanismos para o processamento das entradas geográficas e climáticas, o controle da escala temporal, a configuração de parâmetros e o registro das informações produzidas durante as execuções.

A implementação deve manter correspondência com a especificação conceitual do modelo. Quando uma decisão técnica exigir uma adaptação da representação inicialmente definida, essa alteração deve ser documentada para preservar a rastreabilidade entre o modelo e o código.

A arquitetura do projeto e os procedimentos técnicos utilizados na implementação são apresentados detalhadamente no Capítulo 5.

\section{Procedimentos de verificação}

A verificação tem como objetivo analisar se a implementação apresenta comportamento consistente com as regras e condições definidas para o modelo. Esse procedimento não corresponde, isoladamente, à validação epidemiológica dos resultados.

Os procedimentos de verificação são organizados em diferentes aspectos do simulador, incluindo os estados dos agentes, as transições epidemiológicas, as condições de inicialização, o controle temporal, os parâmetros e os indicadores produzidos.

\subsection{Verificação dos estados e transições}

As transições entre estados devem ser analisadas de acordo com a especificação do modelo. Para isso, devem ser verificadas as condições necessárias para cada transição e a impossibilidade de ocorrência de estados incompatíveis com as regras definidas.

No caso dos indicadores epidemiológicos, também devem ser observadas relações de consistência entre as contagens registradas. Quando os estados forem mutuamente exclusivos e representarem a totalidade de uma população, as somas dos indicadores devem corresponder à quantidade de agentes considerada.

\subsection{Verificação das condições iniciais e dos parâmetros}

As condições de inicialização devem ser verificadas para garantir que as populações e os componentes do ambiente sejam criados de acordo com as configurações definidas.

Os parâmetros utilizados durante a execução também devem ser analisados quanto à sua interpretação e aos limites estabelecidos. Alterações em um parâmetro devem produzir efeitos compatíveis com a regra à qual ele está associado, respeitadas as características estocásticas do modelo.

\subsection{Verificação temporal}

A escala temporal utilizada pela simulação deve apresentar correspondência clara entre os ciclos computacionais e os períodos representados no modelo.

A verificação temporal inclui a análise da frequência de atualização dos processos, da leitura dos dados ambientais e do intervalo utilizado para o registro das saídas. Essa etapa busca evitar interpretações incorretas decorrentes de diferenças entre o tempo computacional e o período representado.

\subsection{Verificação das saídas}

Os indicadores produzidos pelo simulador devem ser analisados quanto à consistência de seus valores e à interpretação das informações registradas.

Entre as verificações previstas estão:

\begin{itemize}
    \item ausência de contagens negativas;
    \item consistência entre o total de agentes e os indicadores de estado;
    \item identificação da escala temporal associada aos registros;
    \item correspondência entre os parâmetros utilizados e os resultados registrados;
    \item análise da possibilidade de reprodução das execuções sob condições equivalentes.
\end{itemize}

Os resultados dos procedimentos de verificação são apresentados posteriormente no capítulo de resultados e avaliação.

\section{Protocolo experimental}

Os experimentos computacionais são utilizados para analisar o comportamento do modelo sob condições previamente definidas. Para garantir a interpretação dos resultados, cada experimento deve registrar as configurações utilizadas, os parâmetros alterados e as condições de execução.

O protocolo experimental considera a definição de um cenário de referência e, quando aplicável, a comparação com outros cenários obtidos pela alteração controlada de parâmetros ou condições do modelo.

\subsection{Cenário de referência}

O cenário de referência corresponde à configuração utilizada como base para a análise do comportamento do simulador. Seus parâmetros e condições iniciais devem ser documentados para permitir a interpretação e a reprodução das execuções.

O cenário de referência não deve ser interpretado automaticamente como uma reprodução exata da situação epidemiológica observada no município. Seu propósito é fornecer uma configuração definida para a realização dos experimentos.

\subsection{Sementes e replicações}

Como o simulador possui componentes estocásticos, diferentes execuções podem produzir resultados distintos mesmo quando utilizam os mesmos parâmetros.

Para analisar essa variabilidade, os experimentos devem considerar múltiplas replicações. O controle das sementes pseudoaleatórias permite identificar as condições utilizadas em cada execução e contribui para a reprodutibilidade dos resultados.

A utilização de uma mesma semente, juntamente com parâmetros e condições equivalentes, deve ser utilizada como procedimento de verificação da reprodutibilidade, quando tecnicamente aplicável.

\subsection{Comparação entre cenários}

A comparação entre cenários deve ser realizada mediante alterações explícitas nas configurações do experimento. Os parâmetros modificados e os indicadores utilizados na comparação devem ser registrados para evitar que diferenças nos resultados sejam atribuídas a fatores não documentados.

A escolha dos cenários deve estar relacionada aos objetivos da pesquisa e às hipóteses que se pretende investigar. Caso sejam realizadas análises de sensibilidade, os parâmetros selecionados e os intervalos avaliados devem ser apresentados de forma explícita.

\section{Indicadores e análise dos resultados}

A análise dos resultados é realizada a partir dos indicadores produzidos durante as execuções do simulador. Os dados registrados devem permitir a identificação das condições da simulação e do comportamento das populações ao longo do tempo.

Entre os indicadores de interesse podem estar as quantidades de agentes em diferentes estados epidemiológicos, informações relacionadas às populações vetoriais e variáveis ambientais utilizadas durante a execução.

Os indicadores são analisados de acordo com os objetivos de cada experimento. Para modelos estocásticos, a interpretação não deve se limitar a uma única execução. Sempre que forem realizadas múltiplas replicações, a análise deve considerar a variabilidade observada entre os resultados.

A apresentação dos resultados pode utilizar tabelas, gráficos e medidas descritivas adequadas aos indicadores disponíveis. A análise deve distinguir resultados efetivamente obtidos durante as execuções de comportamentos esperados ou apenas previstos pela especificação do modelo.

As contagens produzidas pelo simulador também não devem ser consideradas automaticamente equivalentes a indicadores epidemiológicos oficiais. Uma comparação com dados observados exige a análise da definição de cada indicador, da cobertura espacial e da compatibilidade temporal entre as fontes.

\section{Limitações metodológicas}

A metodologia adotada apresenta limitações decorrentes das simplificações necessárias para a construção do modelo, da disponibilidade dos dados e das condições de experimentação.

A utilização de uma população sintética, a agregação espacial, a representação simplificada de processos biológicos e epidemiológicos e a definição de parâmetros constituem exemplos de decisões que podem influenciar o comportamento da simulação.

Além disso, a qualidade das conclusões obtidas depende das verificações e dos experimentos efetivamente realizados. A análise da estrutura do código ou a inspeção de resultados isolados não substitui a execução de testes sistemáticos.

O modelo também não deve ser interpretado como um sistema de previsão epidemiológica validado sem procedimentos específicos de calibração e comparação com dados observados. Os resultados devem ser analisados dentro dos limites definidos pelo propósito da pesquisa e pelas hipóteses adotadas.

A explicitação dessas limitações busca evitar extrapolações que ultrapassem o escopo do estudo e fornecer uma base para futuras melhorias na representação do modelo, nos dados utilizados e nos procedimentos de avaliação.

